% =============================================================================
%  البرنامج السنوي — المعلوماتية — جذع مشترك علوم وتكنولوجيا
%  السنة الدراسية: 2026 / 2027
% =============================================================================
%  البناء (من هذا المجلد):
%    xelatex -interaction=nonstopmode -halt-on-error -output-directory=build annual-program-science.tex
% =============================================================================

\documentclass[11pt, a4paper]{article}

% ── المشترك (مسار shared نسبي لمجلد documents/) ──────────────────────────────
\def\shareddir{../shared/}
\input{"../shared/common.tex"}

% ── بدون تذييل (الترويسة فقط) ───────────────────────────────────────────────
\pagestyle{empty}

% ── قائمة الوحدات وعناوين المجالات تأتي من lesson-macros.tex (المشترك) ──────

% ══════════════════════════════════════════════════════════════════════════════
\begin{document}

\officialheader

\begin{center}
  {\large\bfseries\color{primary} البرنامج السنوي لمادة المعلوماتية — السنة الأولى ثانوي  جذع مشترك علوم وتكنولوجيا}
\end{center}

\vspace{3pt}

% ── المجال الأول ────────────────────────────────────────────────────────────
\domain{الأول}{بيئة التعامل مع الحاسوب}
\begin{units}
  \item تقنية المعلومات (IT)
  \item تجميع الحاسوب
  \item نظام التشغيل
  \item لوحة التحكم
  \item حماية الحاسوب
  \item الشبكة المحلية (LAN)
\end{units}

% ── المجال الثاني ───────────────────────────────────────────────────────────
\domain{ الثاني}{الخوارزميات والمخططات الانسيابية}
\begin{units}
  \item المخططات الانسيابية
  \item انشاء المخططات الانسيابية 
  \item مدخل الى الخوارزميات
  \item التعليمات الأساسية 
\end{units}

% ── المجال الثالث ───────────────────────────────────────────────────────────
\domain{الثالث}{المكتبية}
\begin{units}
  \item معالج النصوص 1
  \item معالج النصوص 2 — دمج المراسلات
  \item جداول البيانات 1 
  \item جداول البيانات 2 — الفرز والترتيب
  \item العروض التقديمية
  \item العروض التقديمية 2 — الحركة بين الشراىح 

\end{units}

% ── المجال الرابع ───────────────────────────────────────────────────────────
\domain{الرابع}{تقنيات الويب}
\begin{units}
  \item المتصفح ومحركات البحث
  \item إنشاء صفحة ويب (HTML)
  \item البريد الالكتروني 
  \item استغلال ادوات الاتصال 

\end{units}

\end{document}
